\subsection{Planning two shared prefix queries}
\label{sec:shared-planner-algorithm}
We plan one pair of readings for all requested contrasts.  The inputs are
$m$ normalized coefficient rows $c_r\in\mathbb R^n$, each with range one,
an exact total $E>0$, and all internal prefix queries
$P_kx=\sum_{h\leq k}x_h$, where $x=y/E\geq0$ and $\sum_hx_h=1$.
We assume $n\geq3$ and a budget of two distinct queries.  Define the range
over a consecutive segment by
$\Delta_r(a{:}b)=\max_{a\leq h\leq b}c_{rh}-\min_{a\leq h\leq b}c_{rh}$.

Suppose the first query is $k$ and its answer is $t$.  For a proposed
second query $j\ne k$, the conditional worst width of contrast $r$ is
$\alpha_{kjr}t+\beta_{kjr}(1-t)$, where
\[
 (\alpha_{kjr},\beta_{kjr})=
 \begin{cases}
 (\max\{\Delta_r(1{:}j),\Delta_r(j+1{:}k)\},
             \Delta_r(k+1{:}n)),&j<k,\\
 (\Delta_r(1{:}k),
       \max\{\Delta_r(k+1{:}j),\Delta_r(j+1{:}n)\}),&j>k.
 \end{cases}
\]
Within each side of the first boundary, a worst-case allocation places
its mass in the segment with the larger coefficient range.  Maximizing
over records and over the finite contrast family commutes, so
\[
 g_{kj}(t)=\max_r\{\alpha_{kjr}t+\beta_{kjr}(1-t)\},\qquad
 A=\min_k\max_{t\in[0,1]}\min_{j\ne k}g_{kj}(t)
\]
is the optimal deterministic two-query worst-width guarantee.

\begin{center}
\fbox{\begin{minipage}{0.94\linewidth}
\small
\textbf{Algorithm: two shared queries, exact-total simplex}
\begin{enumerate}
\item For each first boundary $k$, form every affine line
$\ell_{kjr}(t)=(\alpha_{kjr}-\beta_{kjr})t+\beta_{kjr}$.
\item Form $B_k$ from $0$, $1$, and every intersection in $[0,1]$ of two
nonparallel lines, including lines belonging to different second queries.
\item Compute $A_k=\max_{t\in B_k}\min_{j\ne k}\max_r\ell_{kjr}(t)$.
Choose the smallest boundary index attaining $\min_k A_k$.
\item Retrieve that boundary's counter value and calculate its normalized
prefix mass $t$ using the retained endpoint total.
\item Choose the smallest $j\ne k$ attaining $\min_j g_{kj}(t)$,
retrieve it, and calculate all terminal segment-width certificates using
the two shared answers.
\end{enumerate}
\end{minipage}}
\end{center}

Between consecutive points of $B_k$, the ordering of all underlying lines
is fixed.  The maximum within each group and then the minimum over groups
therefore select one affine line on that interval.  Its maximum is at an
endpoint.  This proves that the finite calculation covers every possible
first answer; it is not a grid approximation.  Equal or coincident lines
are handled by the stated index rule.

The mathematical algorithm uses exact comparisons.  The reported
six-interval implementation uses the same construction in floating-point
arithmetic: it treats values within $10^{-12}$ of a minimum as tied and
chooses the smallest index, both offline and online.  Its intersection
routine skips slope differences of magnitude at most $10^{-14}$; a
257-point grid is an additional numerical check.  Exact rational branch
and adversarial certificates support the reported exact assurance claims.
The algorithm above addresses deterministic worst width on the uncapped,
exact-total simplex.  Bounded-error and capacity analyses use the separate
conditional calculations described in their sensitivity protocols;
randomized mean-width and precision-success comparisons retain their own
optimization objectives.
